\documentclass[12pt,a4paper]{article}
\usepackage{fontspec}
\setmainfont{DejaVu Serif}
\usepackage[russian]{babel}
\usepackage{geometry}
\usepackage{amsmath,amssymb,mathtools}
\usepackage{array,booktabs,longtable}
\usepackage{xcolor,hyperref,tikz}
\geometry{left=22mm,right=20mm,top=22mm,bottom=23mm}
\hypersetup{colorlinks=true,linkcolor=blue!55!black}
\setlength{\parindent}{0pt}
\setlength{\parskip}{6pt}
\renewcommand{\arraystretch}{1.3}
\newcommand{\parthead}[1]{\par\medskip\textbf{#1}\par\smallskip}
\begin{document}
\begin{titlepage}
\thispagestyle{empty}
\vspace*{0.29\textheight}
\begin{center}
{\LARGE\bfseries Ревью домашней работы\par}
\vspace{20mm}
{\large Автор: Лёвин Артём Александрович\par}
\vspace{9mm}
{\large 15 октября 2026 года\par}
\end{center}
\vfill
\begin{center}
\begin{tikzpicture}
\draw[blue!30!black,line width=.65pt]
(-3,0) .. controls (-1.8,.5) and (-.9,-.4) .. (0,0)
.. controls (.9,.4) and (1.8,-.5) .. (3,0);
\end{tikzpicture}
\end{center}
\end{titlepage}

\newpage
\section*{Сводная таблица ревью}
\footnotesize
\begin{longtable}{@{}p{1.15cm}p{2.5cm}p{4.2cm}p{2.8cm}p{3.9cm}@{}}
\toprule
\textbf{№} & \textbf{Тема} & \textbf{Суть ошибки} & \textbf{Ваш ответ} & \textbf{Правильный ответ}\\
\midrule
\endfirsthead
\toprule
\textbf{№} & \textbf{Тема} & \textbf{Суть ошибки} & \textbf{Ваш ответ} & \textbf{Правильный ответ}\\
\midrule
\endhead
\hyperref[task:3a]{3а} & Кубическое уравнение & Перебор возможных корней не завершён; дробные кандидаты не проверены. & Не указан & $-\frac13,\frac12,4$\\ \midrule
\hyperref[task:3b]{3б} & Группировка многочлена & Решение обязательного подпункта отсутствует. & Не указан & $-3$\\ \midrule
\hyperref[task:4a]{4а} & Метод интервалов & Потерян знак минус при решении $2x+3=0$; неверные промежутки. & Промежутки записаны неверно & $\left[-\frac32;-\sqrt2\right]\cup[1;\sqrt2]\cup[7;+\infty)$\\ \midrule
\hyperref[task:4b]{4б} & Знаки произведения & Пропущена точка $2$, не учтён двойной корень $-1$, включены нули. & $(-\infty;-2]\cup\{-1\}\cup[0;1]$ & $(-\infty;-2)\cup(0;2)$\\ \midrule
\hyperref[task:4c]{4в} & Рациональное неравенство & Включены запрещённые значения знаменателя и лишний луч. & $[-1;0]\cup\{1\}\cup(3;+\infty)\cup\{3\}$ & $[-1;0)$\\
\bottomrule
\end{longtable}
\normalsize

\newpage
\thispagestyle{empty}
\vspace*{\fill}
\begin{center}{\Huge\bfseries Разбор ошибок}\end{center}
\vspace*{\fill}

\newpage
\phantomsection\label{task:3a}
\section*{Задание 3а}
\parthead{Текст задания}
Решите уравнение:
\[
6x^3-25x^2+3x+4=0.
\]
\parthead{Ваше решение}
Вы выписали целые делители свободного члена: $\pm1$, $\pm2$, $\pm4$. Проверили значения $x=-1$ и $x=1$ и получили ненулевые значения многочлена. Начали проверять $x=2$, но не закончили вычисление. Решение на этом оборвалось.
\parthead{Ваш ответ}
Не указан.
\parthead{Ошибка}
\textcolor{red}{Ошибка:} проверка возможных корней не завершена, а рассмотрены только целые делители свободного члена.
\parthead{Где именно возникает ошибка}
Последний законченный правильный шаг: $P(1)=6-25+3+4=-12\ne0$. Также верно $P(-1)=-30\ne0$. Первый пропуск --- незавершённая проверка $x=2$ и отсутствие проверки $x=4$ и дробных кандидатов.
\parthead{Почему это неверно}
При целых коэффициентах рациональный корень в виде несократимой дроби $p/q$ имеет числитель, делящий свободный член, и знаменатель, делящий старший коэффициент. В данном уравнении свободный член равен $4$, а старший коэффициент равен $6$, поэтому возможны и дробные корни. Проверка только нескольких целых значений не означает, что уравнение решено. После нахождения одного корня надо разделить многочлен на соответствующий линейный множитель и найти остальные.
\parthead{\textcolor{blue!60!black}{Ключевая идея:}}
Найти один рациональный корень и разложить кубический многочлен на линейный и квадратный множители.
\parthead{Исправление от последнего правильного шага}
Доведём начатую проверку до конца:
\[
P(2)=6\cdot8-25\cdot4+3\cdot2+4=48-100+6+4=-42\ne0.
\]
Проверим $4$:
\[
P(4)=6\cdot64-25\cdot16+12+4=384-400+16=0.
\]
Таким образом, $x=4$ является корнем, а $x-4$ --- множителем исходного многочлена.
\parthead{Полное правильное решение}
Выполним деление многочленов:
\[
6x^3-25x^2+3x+4=(x-4)(6x^2-x-1).
\]
Проверка раскрытием скобок даёт исходный многочлен. Поэтому
\[
(x-4)(6x^2-x-1)=0.
\]
Первый множитель даёт $x=4$. Второй обращается в нуль при
\[
6x^2-x-1=0,\qquad D=(-1)^2-4\cdot6\cdot(-1)=25.
\]
Используя формулу корней квадратного уравнения, находим
\[
x=\frac{1\pm\sqrt{25}}{12}=\frac{1\pm5}{12},
\qquad x_1=\frac12,\quad x_2=-\frac13.
\]
Следовательно, все три корня найдены.
\parthead{Проверка}
Поскольку
\[
6x^3-25x^2+3x+4=(x-4)(2x-1)(3x+1),
\]
каждое найденное число обращает один линейный множитель в нуль. Других корней произведение трёх линейных множителей не имеет.
\parthead{\textcolor{teal!60!black}{Итог:}}
\[
\boxed{x\in\left\{-\frac13;\frac12;4\right\}.}
\]
\parthead{Задача на закрепление}
Решите уравнение $6x^3-7x^2-14x+8=0$.

\newpage
\phantomsection\label{task:3b}
\section*{Задание 3б}
\parthead{Текст задания}
Решите уравнение:
\[
x^5+3x^4+2x^3+6x^2+2x+6=0.
\]
\parthead{Ваше решение}
На представленных страницах решение подпункта отсутствует.
\parthead{Ваш ответ}
Не указан.
\parthead{Ошибка}
\textcolor{red}{Ошибка:} уравнение не решено.
\parthead{Где именно возникает ошибка}
Неверного преобразования здесь нет, поскольку решение не начато. Не выполнен первый содержательный шаг --- группировка слагаемых, позволяющая вынести общий множитель.
\parthead{Почему это неверно}
Хотя степень уравнения равна пяти, структура коэффициентов позволяет решить его школьным способом. Первые два слагаемых, следующие два и последние два имеют общие множители, после вынесения которых возникает одинаковая скобка. Поэтому не требуется искать пять корней подбором или применять методы вне школьной программы.
\parthead{\textcolor{blue!60!black}{Ключевая идея:}}
Группировать члены так, чтобы из каждой группы выносилась одна скобка $x+3$.
\parthead{Исправление}
Начнём с равносильного разбиения многочлена:
\[
(x^5+3x^4)+(2x^3+6x^2)+(2x+6)=0.
\]
\parthead{Полное правильное решение}
Последовательно выносим общий множитель из каждой группы:
\begin{align*}
x^4(x+3)+2x^2(x+3)+2(x+3)&=0,\\
(x+3)(x^4+2x^2+2)&=0.
\end{align*}
Произведение равно нулю тогда и только тогда, когда один из множителей равен нулю. Из $x+3=0$ получаем $x=-3$.

Для второго множителя:
\[
x^4+2x^2+2=(x^2+1)^2+1.
\]
Поскольку квадрат действительного числа неотрицателен, правая часть строго больше нуля при любом действительном $x$. Значит, равенство $x^4+2x^2+2=0$ невозможно, и иных действительных корней нет.
\parthead{Проверка}
Подстановка $x=-3$ в разложенное выражение $(x+3)(x^4+2x^2+2)$ даёт нуль. Второй множитель строго положителен, поэтому ответ полный.
\parthead{\textcolor{teal!60!black}{Итог:}}
\[
\boxed{x=-3.}
\]
\parthead{Задача на закрепление}
Решите уравнение $x^5-2x^4+3x^3-6x^2+4x-8=0$.

\newpage
\phantomsection\label{task:4a}
\section*{Задание 4а}
\parthead{Текст задания}
Решите неравенство:
\[
(x^2-8x+7)(x^2-2)(2x+3)\ge0.
\]
\parthead{Ваше решение}
Вы правильно нашли корни $x^2-8x+7$: $1$ и $7$. Также правильно получили корни $x^2-2$: $-\sqrt2$ и $\sqrt2$. Но для третьего множителя записали
\[
2x+3=0\Rightarrow2x=3\Rightarrow x=1{,}5.
\]
Далее получена схема знаков с неправильными критическими точками и записаны неверные промежутки.
\parthead{Ваш ответ}
Указаны, в частности, промежутки $(-\infty;-\sqrt2]$ и $[7;+\infty)$; объединение в целом неверно.
\parthead{Ошибка}
\textcolor{red}{Ошибка:} при переносе числа $3$ потерян знак минус; далее неверно применён метод интервалов.
\parthead{Где именно возникает ошибка}
Последний верный шаг: $2x+3=0$. Первый неверный шаг: $2x=3$. Нужно было записать $2x=-3$, поэтому $x=-\frac32$, а не $1{,}5$. Корни первых двух множителей сохраняем.
\parthead{Почему это неверно}
Перенос слагаемого из одной части уравнения в другую означает прибавление противоположного числа к обеим частям. При вычитании $3$ получаем отрицательное число $-3$. В методе интервалов знаки меняются именно в корнях множителей. Если корень пропущен, а добавлена несуществующая критическая точка, построенная схема не показывает знак исходного произведения.
\parthead{\textcolor{blue!60!black}{Ключевая идея:}}
Правильно найти и упорядочить все корни множителей, определить знак на одном промежутке и затем учитывать кратность корней.
\parthead{Исправление от последнего правильного шага}
\[
2x+3=0\Longleftrightarrow2x=-3\Longleftrightarrow x=-\frac32.
\]
Общий порядок пяти критических точек:
\[
-\frac32<-\sqrt2<1<\sqrt2<7.
\]
В частности, $\frac32>\sqrt2$, поскольку $\frac94>2$; для отрицательных чисел порядок обратный.
\parthead{Полное правильное решение}
Разложим на множители:
\[
x^2-8x+7=(x-1)(x-7),\quad
x^2-2=(x-\sqrt2)(x+\sqrt2).
\]
Исследуем произведение
\[
(x-1)(x-7)(x-\sqrt2)(x+\sqrt2)(2x+3).
\]
При $x>7$ все множители положительны, поэтому произведение положительно. Все пять корней имеют кратность один, так что при переходе через каждый из них знак меняется:
\[
\begin{array}{c|cccccc}
\text{Промежуток}&(-\infty;-\tfrac32)&(-\tfrac32;-\sqrt2)&(-\sqrt2;1)&(1;\sqrt2)&(\sqrt2;7)&(7;+\infty)\\\hline
\text{Знак}&-&+&-&+&-&+
\end{array}
\]
Так как неравенство нестрогое, включаем точки, при которых произведение равно нулю. Берём второй, четвёртый и шестой промежутки вместе с конечными точками.
\parthead{Проверка}
При $x=0$ исходное произведение равно $7\cdot(-2)\cdot3=-42<0$, поэтому промежуток $(-\sqrt2;1)$ не подходит. При $x=8$ все множители положительны, значит, число подходит. В каждой из пяти критических точек произведение равно нулю и условие $\ge0$ выполнено.
\parthead{\textcolor{teal!60!black}{Итог:}}
\[
\boxed{x\in\left[-\frac32;-\sqrt2\right]\cup[1;\sqrt2]\cup[7;+\infty).}
\]
\parthead{Задача на закрепление}
Решите неравенство $(x^2-5x+6)(x^2-3)(2x+1)\ge0$.

\newpage
\phantomsection\label{task:4b}
\section*{Задание 4б}
\parthead{Текст задания}
Решите неравенство:
\[
(x^4-16)(x^3+1)(x^2+x)<0.
\]
\parthead{Ваше решение}
Вы правильно вычислили корни отдельных множителей: $-2$ и $2$; $-1$; $0$ и $-1$. В схеме знаков, однако, появилась точка $1$ вместо $2$, а при переходе через $-1$ знак изменяется. В ответе оказались нули произведения.
\parthead{Ваш ответ}
$(-\infty;-2]\cup\{-1\}\cup[0;1]$.
\parthead{Ошибка}
\textcolor{red}{Ошибка:} неверно составлена схема знаков, проигнорирована чётная кратность корня $-1$, а концы промежутков включены при строгом неравенстве.
\parthead{Где именно возникает ошибка}
Последний правильный этап --- нахождение корней каждого множителя. Первый неверный шаг --- нанесение на числовую прямую $1$ вместо $2$. Кроме того, точка $-1$ даёт нуль сразу двух множителей, поэтому при переходе через неё знак не должен меняться.
\parthead{Почему это неверно}
В неравенстве со знаком $<0$ значения, обращающие произведение в нуль, всегда исключаются. При $x=-1$ оба множителя $x^3+1$ и $x^2+x$ обращаются в нуль. После разложения получается квадрат $(x+1)^2$, который неотрицателен и по обе стороны от $-1$ положителен. Поэтому здесь нет смены знака. Напротив, $x=2$ --- реальный корень первого множителя, который обязательно нужно учитывать.
\parthead{\textcolor{blue!60!black}{Ключевая идея:}}
Следить не только за расположением корней, но и за их кратностью: при чётной кратности знак не меняется.
\parthead{Исправление от последнего правильного шага}
На прямой должны стоять $-2<-1<0<2$. В точке $-1$ корень повторяется дважды. Никакой критической точки $1$ у исходного произведения нет.
\parthead{Полное правильное решение}
По формулам разности квадратов и суммы кубов:
\begin{align*}
x^4-16&=(x-2)(x+2)(x^2+4),\\
x^3+1&=(x+1)(x^2-x+1),\\
x^2+x&=x(x+1).
\end{align*}
Поэтому всё произведение равно
\[
x(x-2)(x+2)(x+1)^2(x^2+4)(x^2-x+1).
\]
Здесь $x^2+4>0$ всегда, а
\[
x^2-x+1=\left(x-\frac12\right)^2+\frac34>0
\]
для любого действительного $x$. Квадрат $(x+1)^2$ строго положителен при $x\ne-1$ и равен нулю в точке $-1$. За исключением нулей произведения его знак совпадает со знаком $x(x-2)(x+2)$.

Справа от $2$ этот продукт положителен. При переходе через простой корень $2$ знак меняется на отрицательный. После простого корня $0$ становится положительным, после $-2$ --- отрицательным:
\[
\begin{array}{c|cccc}
\text{Промежуток}&(-\infty;-2)&(-2;0)&(0;2)&(2;+\infty)\\\hline
\text{Знак}&-&+&-&+
\end{array}
\]
Точка $-1$ лежит внутри положительного промежутка $(-2;0)$ и обращает произведение в нуль, но не меняет знак. Поэтому она не добавляет решения.
\parthead{Проверка}
При $x=-3$ знаки исходных трёх множителей: положительный, отрицательный, положительный, значит, произведение отрицательно. При $x=1$ знаки: отрицательный, положительный, положительный, значит, произведение отрицательно. При $x=-\frac12$ знаки: отрицательный, положительный, отрицательный, значит, произведение положительно. Нули $-2,-1,0,2$ исключены строгим знаком.
\parthead{\textcolor{teal!60!black}{Итог:}}
\[
\boxed{x\in(-\infty;-2)\cup(0;2).}
\]
\parthead{Задача на закрепление}
Решите неравенство $(x^2-9)(x^3+8)(x^2+2x)<0$.

\newpage
\phantomsection\label{task:4c}
\section*{Задание 4в}
\parthead{Текст задания}
Решите неравенство:
\[
\frac{x^2-1}{(x^2-x)|3-x|}\le0.
\]
\parthead{Ваше решение}
Вы правильно выделили нули числителя $-1$ и $1$ и нули знаменателя $0$, $1$, $3$. Затем при составлении ответа включили запрещённые точки и луч справа от $3$.
\parthead{Ваш ответ}
$[-1;0]\cup\{1\}\cup(3;+\infty)\cup\{3\}$.
\parthead{Ошибка}
\textcolor{red}{Ошибка:} не соблюдена область допустимых значений дроби; неверно определён знак при $x>3$.
\parthead{Где именно возникает ошибка}
Последний правильный шаг --- нахождение нулей знаменателя $x=0,1,3$. Первый неверный шаг --- включение этих точек в ответ. Кроме того, при $x>3$ числитель и все множители знаменателя положительны, так что положительный правый луч не должен входить в решение.
\parthead{Почему это неверно}
Дробь определена лишь при ненулевом знаменателе. Ни равенство нулю числителя, ни сокращение общих множителей не могут сделать разрешённой точку, где знаменатель исходной дроби равен нулю. Здесь $|3-x|$ положителен по обе стороны от $3$ и не создаёт смены знака; при $3$ он равен нулю и само выражение не определено.
\parthead{\textcolor{blue!60!black}{Ключевая идея:}}
Сначала установить ОДЗ, затем сократить общий множитель, не забывая исключённые точки, и исследовать знак оставшегося отношения.
\parthead{Исправление от последнего правильного шага}
Так как
\[
(x^2-x)|3-x|=x(x-1)|3-x|,
\]
исключаем $x=0$, $x=1$, $x=3$. Точка $x=-1$ разрешена: в ней числитель равен нулю, а знаменатель не равен нулю.
\parthead{Полное правильное решение}
Разложим числитель и знаменатель:
\[
\frac{(x-1)(x+1)}{x(x-1)|3-x|}\le0,
\qquad x\ne0,1,3.
\]
При $x\ne1$ можно сократить ненулевой множитель $x-1$:
\[
\frac{x+1}{x|3-x|}\le0,
\qquad x\ne0,1,3.
\]
На ОДЗ модуль $|3-x|$ строго положителен. Поэтому знак совпадает со знаком отношения $\frac{x+1}{x}$.

Если $x<-1$, и числитель, и знаменатель отрицательны, значит, отношение положительно. При $x=-1$ отношение равно нулю и точка разрешена. При $-1<x<0$ числитель положителен, знаменатель отрицателен, отношение отрицательно. При $x>0$ обе величины положительны, поэтому отношение положительно во всех допустимых точках этой части прямой. Следовательно, ответ составляют точка $-1$ и промежуток от $-1$ до $0$ без правого конца.
\parthead{Проверка}
При $x=-\frac12$:
\[
\frac{x^2-1}{(x^2-x)|3-x|}
=\frac{-\frac34}{\frac34\cdot\frac72}=-\frac27<0.
\]
При $x=4$ исходное выражение равно $\frac{15}{12\cdot1}=\frac54>0$. Значения $0$, $1$, $3$ запрещены из-за нулевого знаменателя.
\parthead{\textcolor{teal!60!black}{Итог:}}
\[
\boxed{x\in[-1;0).}
\]
\parthead{Задача на закрепление}
Решите неравенство $\displaystyle\frac{x^2-4}{(x^2-2x)|5-x|}\le0$.
\vfill
\begin{center}
\small Лёвин Артём Александрович эксклюзивно для Ксении Клыковой
\end{center}
\end{document}
